\section{Discussion}
\label{sec:discussion}

The results distinguish three kinds of information obtainable from
quadratic aggregation. Under asymptotic hyperplane convexity, a single
nonconstant globally convex aggregation detects whether a nonempty
strict system has a proper hull. Complete hull descriptions require a
larger class of good aggregations. Their size depends on geometry that
the certificate alone does not record: homogeneous PDLC for three rows
gives a sharp four-bound, whereas HHC without that definiteness can require
uncountably many strict quadratic inequalities even in four variables.
The latter example still has a small exact semidefinite lift and
quantitatively controlled finite approximations.

Several distinctions persist at the boundaries of these results.
Strict feasibility supports the certificate theorem and the comparison
with the Shor projection; closed systems without a strictly feasible
point need separate treatment. Closure also changes representation
cardinality: the explicit open hull needs uncountably many strict
quadratics, while a countable dense good family suffices for its closed
hull. Under PDLC the limiting negative components must retain their
orientation to obtain a valid closed conic description. None of these
operations can be replaced in general by changing inequality signs.

The statements are structural rather than complexity guarantees.
No general procedure for recognizing HHC or for efficiently selecting the
four PDLC multipliers is established here. The exact SDP characterizations
use exact optimal values, and the rational approximation families control
coefficient size without analyzing numerical conditioning. For a
three-dimensional span with many extreme rays, the directional facet
bound improves the known count under its stated cone condition; the
counterexample in Appendix~\ref{app:three-span} rules out a general
constant bound from the complementary condition alone. These limits
define the reach of the proved results without being premises on which
any theorem in the paper remains contingent.
